\subsection{Энергия пружинок}
%%%%%%%%%%%%%%%%%%%%%%%%%%%%%%%%%%%%%%%%%%%%%%%%%%%%
%%%%%%%%%%%%%%%%%%%%%%%%%%%%%%%%%%%%%%%%%%%%%%%%%%%% 
\z[{\begin{wrapfigure}[5]{r}{0.2\textwidth}%
		\vspace{-10pt}%
		\begin{flushright}%
 			\includegraphics[width=0.2\textwidth]{7.6.1}%
		\end{flushright}%
	\end{wrapfigure}%
	}]%
	{Тело массы $m$ падает с высоты $h$ на стоящую вертикально на полу пружину жесткости $k$ и длины $l$. Определите максимальную силу давления на пол.}
%
{$F = mg\left(1 + \sqrt{1 + \frac{2k(h - l)}{mg}}\right)$}
%%%%%%%%%%%%%%%%%%%%%%%%%%%%%%%%%%%%%%%%%%%%%%%%%%%%	
\z{На пружинке жесткости $k$ покоится грузик массой $m$. С высоты $H$ над грузиком отпускают пластилиновый шарик массы $2m$. На сколько максимум опустится грузик?}
%
{$\Delta l = \frac{2mg + \sqrt{4m^2 g^2 + \frac{8mgkH}{3}}}{k}$}
%%%%%%%%%%%%%%%%%%%%%%%%%%%%%%%%%%%%%%%%%%%%%%%%%%%% 
\z[{\begin{wrapfigure}[5]{r}{0.3\textwidth}%
		\vspace{-10pt}%
		\begin{flushright}%
 			\includegraphics[width=0.3\textwidth]{7.6.3}%
		\end{flushright}%
	\end{wrapfigure}%
	}]%
	{Между шариками массами $m$ и $M$, связанными нитью, вставлена лёгкая пружина жёсткостью $k$, сжатая на некоторую величину. Система шариков движется со скоростью $v_0$ вдоль прямой, проходящей через центры шариков (см. рисунок). Нить пережигают, и один из шариков останавливается. Найти начальную величину сжатия пружины.}
%
{$x = v_0 \sqrt{\frac{m}{k}\left(1 + \frac{m}{M}\right)}$}
%%%%%%%%%%%%%%%%%%%%%%%%%%%%%%%%%%%%%%%%%%%%%%%%%%%%
\z[{\begin{wrapfigure}[7]{r}{0.15\textwidth}%
		\vspace{-10pt}%
		\begin{flushright}%
 			\includegraphics[width=0.15\textwidth]{7.6.4}%
		\end{flushright}%
	\end{wrapfigure}%
	}]%
	{Альпинистская капроновая верёвка подчиняется закону Гука, пока не разрывается при силе натяжения $T = 22000\,\text{Н}$, будучи растянутой на $\alpha = 25\%$ от своей первоначальной длины. Стандартный способ испытания верёвки такой: один конец верёвки длиной $L$ закрепляют на стене, и с высоты, равной $L$, сбрасывают груз массой $m$, привязанный к другому концу (см. рисунок). При каком максимальном грузе $m$ верёвка обязана выдержать рывок?}
%
{$m = \frac{T \alpha}{2g(2 + \alpha)} \approx 122\,\text{кг}$}
%%%%%%%%%%%%%%%%%%%%%%%%%%%%%%%%%%%%%%%%%%%%%%%%%%%%	
\z[{\begin{wrapfigure}[6]{r}{0.3\textwidth}%
		\vspace{+10pt}%
		\begin{flushright}%
 			\includegraphics[width=0.3\textwidth]{7.6.5}%
		\end{flushright}%
	\end{wrapfigure}%
	}]%
	{Грузы с одинаковыми массами связаны нитью. Между ними вставлена лёгкая пружина жёсткостью $k$, сжатая на величину $x_0$. Грузы движутся со скоростью $v$ вдоль прямой, составляющей угол $\alpha$ с осью системы (см. рисунок). После пережигания нити один из грузов полетел перпендикулярно первоначальному направлению движения. Найти массу груза.}
%
{$m = \frac{k x_0^2 \cos^2 \alpha}{2v^2}$}